\chapter{集合论基础与选择原理}
\label{chap:appendix-set-theory}

\begin{chapterguide}
分析学通常把数、函数和空间作为已经存在的对象使用。本附录说明这些对象在 ZFC 中为何能够形成集合，并把后文偶尔涉及的选择原则分开。这里不以集合论本身为研究对象；所保留的细节以消除存在性疑问、核对选择强度为限。
\end{chapterguide}

\section{一阶语言、经典逻辑与集合论的论域}

集合论的一阶语言只有等号和二元关系符号 \(\in\) 两个非逻辑符号。变量遍历集合，原子公式为 \(x=y\) 与 \(x\in y\)，一般公式由逻辑联结词和量词有限次组成。公式中的变量可能是自由的，也可能被量词约束；公理是没有自由变量的句子，或由一个公式按参数统一生成的公理模式。

“任意集合”“全体集合”是对变量论域的说明，并不表示存在一个由所有集合组成的集合。带参数的记号
\[
 \{x\in A:\varphi(x,p_1,\ldots,p_k)\}
\]
只有在母集 \(A\) 已经给定时才表示集合。第 \ref{chap:mathematical-language} 章采用的受限概括，正是这种用法。

本书使用经典逻辑，因而可用排中律、反证法和双重否定消去。集合论公理只保证对象存在；对象上的运算能否由代表元给出、所得定义是否与选择无关，仍须另作数学证明。例如商集的存在性与商集上运算的良定义性是两件不同的事。

\section{Zermelo--Fraenkel 公理体系}

ZF 包含下列公理和公理模式。所列形式足以说明本书中的集合构造；不同文献可能把空集公理从其他公理推出，或采用等价的表述。
\begin{description}[style=nextline,leftmargin=3em]
 \item[外延公理] 若两个集合有相同元素，则二者相等。
 \item[空集公理] 存在没有元素的集合 \(\varnothing\)。
 \item[配对公理] 对任意集合 \(x,y\)，存在集合 \(\{x,y\}\)。
 \item[并集公理] 对任意集合 \(A\)，存在 \(\bigcup A=\{x:\exists B\in A,\ x\in B\}\)。
 \item[幂集公理] 对任意集合 \(A\)，存在由 \(A\) 的全部子集组成的 \(\mathcal P(A)\)。
 \item[无穷公理] 存在一个集合 \(I\)，使 \(\varnothing\in I\)，并且 \(x\in I\) 蕴含 \(x\cup\{x\}\in I\)。
 \item[分离公理模式] 对每个公式 \(\varphi\) 和集合 \(A\)，集合 \(\{x\in A:\varphi(x)\}\) 存在。公式可以含有已经给定的集合参数。
 \item[替代公理模式] 若公式 \(\varphi(x,y)\) 对每个 \(x\in A\) 唯一确定一个 \(y\)，则这些 \(y\) 组成集合。
 \item[正则公理] 每个非空集合 \(A\) 含有元素 \(x\)，使 \(x\cap A=\varnothing\)。
\end{description}

ZF 加上选择公理称为 ZFC。配对与并集给出单点集、有序有限组和有限并；幂集提供关系、函数图、拓扑和 \(\sigma\)-代数所在的母集；无穷公理保证归纳集存在。所有归纳子集之交给出 von Neumann 自然数集 \(\omega\)，其中
\[
 0=\varnothing,\qquad n+1=n\cup\{n\}.
\]
它的最小性给出归纳法，递归定理则需要替代公理把各阶段构造的值收集成集合。

正则公理排除 \(x\in x\) 和无限下降成员链。初等分析几乎不直接使用它，但它保证集合可以按累积层级组织。

\section{分离、替代、Russell 悖论与受限概括}

外延公理把集合相等化为双包含。分离模式从一个已有集合中筛选元素；替代模式则把已有集合中的每个元素送到唯一新对象，并把所得像收集成集合。若 \(F\) 是 \(A\) 上的函数式关系，则
\[
 F[A]=\{F(a):a\in A\}
\]
由替代得到，而 \(\{a\in A:F(a)\in B\}\) 由分离得到。

不能删去受限概括中的母集。若允许对任意性质写出集合
\[
 R=\{x:x\notin x\},
\]
便会得到 \(R\in R\Longleftrightarrow R\notin R\)。这个矛盾说明“满足某性质的一切对象”不一定组成集合，并不说明公式 \(x\notin x\) 本身非法。

\begin{proposition}
不存在由所有集合组成的集合。
\end{proposition}

\begin{proof}
若 \(U\) 包含所有集合，分离模式给出
\[
 R_U=\{x\in U:x\notin x\}.
\]
集合 \(R_U\) 也应属于 \(U\)，于是定义给出 \(R_U\in R_U\) 当且仅当 \(R_U\notin R_U\)，矛盾。
\end{proof}

集合族记号也应带有一个集合指标。若 \(I\) 是集合且 \(i\mapsto A_i\) 是定义在 \(I\) 上的函数，则 \(\{A_i:i\in I\}\) 由替代形成集合，随后才可取并、交或幂集。正文中“任意族”均默认满足这一条件。

\section{关系、函数、商结构和有序对编码}

采用 Kuratowski 有序对
\[
 (x,y)=\{\{x\},\{x,y\}\}.
\]
由外延公理直接核对可得
\[
 (x,y)=(x',y')\quad\Longleftrightarrow\quad x=x',\ y=y'.
\]
若 \(x\in X\)、\(y\in Y\)，则 \((x,y)\in\mathcal P(\mathcal P(X\cup Y))\)。因此 \(X\times Y\) 可由分离从这个母集中取得。二元关系是笛卡儿积的子集；函数是满足全定义性与单值性的关系，并与定义域和陪域一同使用。

若 \(\sim\) 是 \(X\) 上的等价关系，则
\[
 [x]=\{y\in X:y\sim x\}
\]
由分离得到，映射 \(x\mapsto[x]\) 的像由替代得到，这个像就是商集 \(X/{\sim}\)。若拟用
\[
 [x]\star[y]=[F(x,y)]
\]
定义运算，还必须证明 \(x\sim x'\)、\(y\sim y'\) 蕴含 \(F(x,y)\sim F(x',y')\)。第 \ref{chap:cauchy-construction} 章的 Cauchy 列商构造正是这一原则的主要实例。

有限序列可用函数编码，无穷序列是定义在 \(\omega\) 或 \(\Npos\) 上的函数。一个递归定义若在每个阶段唯一给出下一项，递归定理产生整条序列，不需要选择公理；若每一步只有“至少存在一个”候选而没有规范规则，才可能需要选择原则。

\section{序数、基数、超限递归与 Hartogs 引理}

集合 \(\alpha\) 称为\emph{传递集}，若 \(x\in y\in\alpha\) 蕴含 \(x\in\alpha\)。序数是由 \(\in\) 良序的传递集。von Neumann 约定下，每个序数恰由所有较小序数组成：
\[
 \alpha=\{\beta:\beta<\alpha\},\qquad \beta<\alpha\Longleftrightarrow\beta\in\alpha.
\]
后继序数为 \(\alpha+1=\alpha\cup\{\alpha\}\)；非零且不是后继的序数称为极限序数。自然数是有限序数，\(\omega\) 是最小无限序数。

\begin{theorem}[超限归纳]
设性质 \(P(\alpha)\) 满足：对每个序数 \(\alpha\)，若 \(P(\beta)\) 对所有 \(\beta<\alpha\) 成立，则 \(P(\alpha)\) 成立。于是 \(P\) 对所有序数成立。
\end{theorem}

\begin{proof}
若失败序数形成非空集合，序数的良序性给出其中最小者 \(\alpha\)。对每个 \(\beta<\alpha\)，\(P(\beta)\) 均成立，归纳假设便推出 \(P(\alpha)\)，矛盾。实际使用时只需在某个预先给定的序数内部取最小反例，因而不会形成“所有序数的集合”。
\end{proof}

\begin{theorem}[超限递归]
若规则 \(G\) 对每个序数 \(\alpha\) 和定义在 \(\alpha\) 上的函数 \(g\) 唯一指定集合 \(G(\alpha,g)\)，则对每个序数 \(\gamma\)，存在唯一函数 \(f\) 定义在 \(\gamma\) 上，满足
\[
 f(\alpha)=G(\alpha,f|_\alpha)\qquad(\alpha<\gamma).
\]
\end{theorem}

证明按 \(\gamma\) 作超限归纳：后继阶段添加唯一新值，极限阶段取此前相容函数之并。替代公理把各阶段的值收集成集合；唯一性由“第一处分歧”反证得到。

两个集合等势，若它们之间存在双射。一个集合的基数是它的良序类型中最小的初始序数；在不预设选择公理时，并非每个集合都已经知道可良序，因此下述引理尤其重要。

\begin{theorem}[Hartogs 引理]
对每个集合 \(X\)，存在序数 \(h(X)\)，使得不存在从 \(h(X)\) 到 \(X\) 的单射。
\end{theorem}

\begin{proof}
所有“\(X\) 的某个子集上的良序关系”组成
\[
 \mathcal W\subseteq\mathcal P(X)\times\mathcal P(X\times X),
\]
所以 \(\mathcal W\) 是集合。把其中同构的良序视为同一类型；这些类型仍组成集合，并按“同构于真初段”排序。良序比较定理说明任意两个类型可比，而且每个非空类型族有最小元，故这些类型的序型是一个序数 \(h(X)\)。

若存在单射 \(j:h(X)\to X\)，便可把 \(h(X)\) 的良序搬到子集 \(j[h(X)]\subseteq X\) 上；其序型应属于构造 \(h(X)\) 时收集的类型集合，因而严格小于 \(h(X)\)。这等于断言序数 \(h(X)\) 严格小于自身，矛盾。
\end{proof}

Hartogs 引理在 ZF 中成立；它不以“每个集合都有基数”为前提。它将“递归不能无限产生互异元素”转化为一个集合论上可核验的终止论证。

\section{选择公理、Zorn 引理、良序原理及等价性}

\begin{definition}[选择公理]
若 \(\mathcal A\) 是由非空集合组成的集合，则存在函数
\[
 c:\mathcal A\to\bigcup\mathcal A,\qquad c(A)\in A.
\]
\end{definition}

\begin{definition}[Zorn 引理]
若偏序集 \(P\) 的每条链在 \(P\) 中都有上界，则 \(P\) 有极大元。
\end{definition}

\begin{definition}[良序原理]
每个集合都能赋予一个良序，即该全序的每个非空子集都有最小元。
\end{definition}

\begin{theorem}
在 ZF 中，选择公理、Zorn 引理和良序原理相互等价。
\end{theorem}

\begin{proof}
先设选择公理成立。对非空集合 \(X\)，令 \(c\) 是 \(\mathcal P(X)\setminus\{\varnothing\}\) 上的选择函数。按序数递归：只要此前所选元素尚未穷尽 \(X\)，就令
\[
 x_\alpha=c\bigl(X\setminus\{x_\beta:\beta<\alpha\}\bigr).
\]
若递归在 \(h(X)\) 以前从不停止，\(\alpha\mapsto x_\alpha\) 就是从 \(h(X)\) 到 \(X\) 的单射，与 Hartogs 引理矛盾。因此某一阶段 \(\gamma<h(X)\) 恰好取尽 \(X\)，用 \(x_\alpha<x_\beta\Longleftrightarrow\alpha<\beta\) 便得到 \(X\) 的良序。故选择公理推出良序原理。

由良序原理立刻得到选择公理：先良序 \(\bigcup\mathcal A\)，再令 \(c(A)\) 为 \(A\) 的最小元。

再由良序原理推出 Zorn 引理。良序 \(P=\{p_\alpha:\alpha<\gamma\}\)，按序递归构造链：考察 \(p_\alpha\) 时，当且仅当它与此前收入链的每个元素可比，才把它加入。所得链 \(C\) 是极大链，因为任何与 \(C\) 全部元素可比的点在被考察时都会收入 \(C\)。设 \(u\) 是 \(C\) 的上界。若有 \(v>u\)，则 \(v\) 与 \(C\) 中每个元素可比，所以 \(v\in C\)；但又有 \(v\le u\)，矛盾。故 \(u\) 是 \(P\) 的极大元。

最后设 Zorn 引理成立。考虑所有二元组 \((Y,<_Y)\)，其中 \(Y\subseteq X\) 且 \(<_Y\) 良序 \(Y\)；若前者是后者的初段并且两个序在该初段上一致，就称前者小于后者。链的并带有相容的并序，仍是一个良序子集，故 Zorn 引理给出极大元 \((Y,<_Y)\)。若 \(Y\ne X\)，取任意 \(x\in X\setminus Y\)，把 \(x\) 接在 \(Y\) 的全部元素之后，得到真扩张，矛盾。因此 \(Y=X\)，良序原理成立。
\end{proof}

“极大元”不等于“最大元”。Zorn 引理只保证某个元素不能再严格增大，不保证它大于偏序集中的一切元素。后文用 Zorn 引理取得 Hamel 基、极大理想或线性泛函延拓时，都必须先证明链上界确实仍属于所考察的偏序集。

\section{Tychonoff 定理的版本与选择强度}

一般 Tychonoff 定理断言任意族紧空间的乘积在乘积拓扑下紧。证明常把开覆盖改写为具有有限交性质的闭集族，或使用超滤子；在 ZF 中，这一一般命题与选择公理等价。

若只要求每个因子为紧 Hausdorff 空间，相应乘积定理可由\emph{超滤子引理}（也称 Boolean 素理想定理的一种形式）推出。该原则在 ZF 中严格弱于完整选择公理，但仍不能由 ZF 证明。对可数个紧度量空间，还可以在适当的可数选择原则下用逐坐标抽取和对角法证明序列版本；这与任意指标族的一般乘积紧致不可混为一谈。

本附录只登记选择强度。乘积拓扑、网、滤子、超滤子和紧致性的完整定义与证明放在一般拓扑各章；在那些概念建立以前，正文不会调用 Tychonoff 定理。这里所说的“等价”均指在 ZF 上增加相应命题后与所列选择原则互相可推。

\section{本书各结论使用的选择原则}

以下几类选择容易混淆。
\begin{itemize}
 \item 从 \(\N\) 的非空子集取最小元、从已固定枚举中取第一个合格对象，是规范选择，不使用选择公理。
 \item 对有限多个非空集合逐一取元可在 ZF 中证明；对可数族非空集合同时取元是可数选择。
 \item 若关系 \(R\) 满足每个状态都有一个 \(R\)-后继，由一个初态递归选择无限轨道通常使用依赖选择。
 \item “可数个至多可数集之并仍至多可数”在没有预先给出各集合枚举时通常使用可数选择；若枚举已经作为数据给出，则对角编排是规范构造。
 \item 任意向量空间有 Hamel 基、一般形式的 Hahn--Banach 定理、每个非零含幺交换环有极大理想，通常经 Zorn 引理使用完整选择公理或与之相关的较弱原则。
 \item 任意乘积的紧致性依赖上一节所区分的 Tychonoff 版本。
\end{itemize}

前三编中的二分区间、最小合格指标和有理数固定枚举均可作规范选择。若证明写成“对每个 \(n\) 任取一个对象”，应检查这些对象能否由最小指标、固定稠密序列或显式公式确定；不能时，正文会注明所用选择原则。全书默认工作于 ZFC，因此注明强度是为了说明证明结构，而不是限制合法性。

\section*{习题}
\addcontentsline{toc}{section}{习题}

\begin{enumerate}[label=\textbf{A.\arabic*.}]
 \exerciseitem{A.1} 用外延公理证明空集唯一，并由配对公理构造单点集。
 \exerciseitem{A.2} 区分下列构造使用分离还是替代：偶数集、函数像 \(f[A]\)、等价类 \([x]\)、商集 \(X/{\sim}\)。
 \exerciseitem{A.3} 从 Kuratowski 编码证明有序对的坐标唯一性。
 \exerciseitem{A.4} 证明不存在全集，并说明该证明没有断言性质 \(x\notin x\) 本身非法。
 \exerciseitem{A.5} 证明一个分割对应的等价关系及一个等价关系对应的分割均形成集合。
 \exerciseitem{A.6} 证明良序原理推出选择公理。
 \projectexerciseitem{A.7} \ 【前置：偏序与超限递归；预计用时：90分钟】沿正文构造补全 Zorn 引理推出良序原理时“链的并仍为良序”的全部细节。
 \optionalexerciseitem{A.8} 分析“可数族至多可数集之并仍至多可数”的证明中选择函数出现的位置。
 \optionalexerciseitem{A.9} 证明若每个成员集合都已指定一个元素，则形成选择函数不再需要选择公理。
 \openexerciseitem{A.10} 查阅一个不能在 ZF 中证明、但在 ZFC 中成立的分析命题，说明其选择强度。
\end{enumerate}
